\subsection{The small resolution as a smooth-side mirror}
\label{subsec:x9-mirror-mutation}

The holomorphic families can also be compared explicitly. The general
ambient fiber of the smoothing of $X_9$ is toric: its fan, reconstructed
in \cite[Section~11]{Wuebben:2026smoothing}, has primitive rays
\begin{equation}
 \begin{aligned}
 r_0&=(-1,-3,-1,0),&r_1&=(-1,-3,0,-1),\\
 r_2&=(-1,-1,0,0),&r_3&=(-1,0,0,0),\\
 r_4&=(0,1,0,0),&r_5&=(1,-1,0,0),&r_6&=(1,3,1,1).
 \end{aligned}
 \label{eq:x9-smooth-mirror-rays}
\end{equation}
Its maximal cones are
$1456,1346,0456,0346,0145,0134,01256,1236,0236,0123$.
This is a Gorenstein Fano fourfold; a general anticanonical
hypersurface is the smooth fiber $X_{9,t}$. Consequently
$P=\operatorname{conv}(r_0,\ldots,r_6)$ is the Newton polytope
of its Batyrev mirror \cite{Batyrev:1994}.

\begin{proposition}\label{prop:x9-mirror-mutation}
For generic $(v,a,b)$, the projective small resolution of the
four-node mirror on $\Sigma$ is birational to a smooth crepant
Batyrev mirror of $X_{9,t}$. The birational map preserves the
residue three-form, and the two smooth projective models are
related by flops. The parameter family $(v,a,b)$ has generically
full rank in the three-dimensional complex-structure moduli space.
\end{proposition}

\begin{proof}
Write \eqref{eq:x9-finite-bulk-factorization} as
$\mathcal F=fg+sk$, with
$f=1+xy$, $g=1+x^{-1}+vy^{-1}+r$ and $k=r(br-1)+as$.
Blow up the compact Weil divisor given in this chart by $(s,f)$.
In the affine blowup chart $s=Tf$, the strict transform is
\begin{equation}
 g+Tr(br-1)+aT^2f=0.
 \label{eq:x9-mutation-blowup-chart}
\end{equation}
On its dense torus set $u=r$, $p=x/r^2$ and $q=y/r^2$.
Dividing \eqref{eq:x9-mutation-blowup-chart} by the torus unit $uT$
gives
\begin{equation}
 \begin{aligned}
 F_{\rm sm}={}&\frac{1}{Tu^3p}+\frac{v}{Tu^3q}
       +\frac{1}{Tu}+\frac{1}{T}+bu
       +\frac{aT}{u}+aTu^3pq-1=0.
 \end{aligned}
 \label{eq:x9-mutated-mirror}
\end{equation}
The exponents of its seven nonconstant monomials, in the order
shown, are precisely \eqref{eq:x9-smooth-mirror-rays}.
The inverse map on the dense open set $1+u^4pq\ne0$ is
\begin{equation}
 x=u^2p,\qquad y=u^2q,\qquad r=u,\qquad
 s=T(1+u^4pq).
 \label{eq:x9-mutation-map}
\end{equation}
Thus this is a birational map from the actual small resolution,
not just a comparison of Newton-polytope invariants.

The three-form is preserved as well. Put $h=s/r$ and
$G=1+u^4pq$. In the integral old torus coordinates $(h,u,p,q)$,
the original Laurent equation is $\mathcal F/(rs)=0$.
The substitution $h=TG/u$ changes this Laurent polynomial
exactly to $F_{\rm sm}$. Moreover
\begin{equation}
 d\log h\wedge d\log u\wedge d\log p\wedge d\log q
 =d\log T\wedge d\log u\wedge d\log p\wedge d\log q,
 \label{eq:x9-mutation-volume}
\end{equation}
because $d\log G$ involves only $u,p,q$.
Taking residues proves the assertion, up to the common orientation
choice. This is a Laurent mutation in the usual volume-preserving
sense \cite[Definition~1.1]{Ilten:2012Mutation}.

We must check that the coefficient restrictions in
\eqref{eq:x9-mutated-mirror} do not impose an extra specialization.
The lattice points of $P$ are its seven vertices, the origin and
$e_*=(0,-1,0,0)$. The point $e_*$ lies in the relative interior
of the facet with inward normal $n_*=(0,1,-2,-2)$.
It pairs nonnegatively with all other primitive facet normals.
The corresponding toric root automorphism acts on anticanonical
monomials by
\begin{equation}
 \chi^m\longmapsto
 \chi^m(1+\lambda/u)^{1+\langle n_*,m\rangle}.
 \label{eq:x9-mirror-root-gauge}
\end{equation}
At the coefficients in \eqref{eq:x9-mutated-mirror}, an added
coefficient $c/u$ transforms to $(c-\lambda+b\lambda^2)/u$;
the constant becomes $-1+2b\lambda$ and the coefficient of
$(Tu)^{-1}$ becomes $1+\lambda$. The derivative of the first
coefficient with respect to $\lambda$ is $-1$ at zero. Thus
$c=0$ is locally a slice for this automorphism. Normalize the
constant coefficient to $-1$. The four torus normalizations
$c_{r_0}=c_{r_2}=c_{r_3}=1$ and $c_{r_6}=c_{r_5}$ have
exponent matrix with rows
$(r_0,r_2,r_3,r_6-r_5)$ and determinant $1$. They leave exactly
$(c_{r_1},c_{r_4},c_{r_5})=(v,b,a)$, with no finite torus
ambiguity.

For completeness, $P$ and its polar have $9$ and $162$ lattice
points; the sums of facet-interior lattice counts are $1$ and
$34$, and both codimension-two corrections vanish. Batyrev's
formulas give $(h^{1,1},h^{2,1})=(123,3)$ and show that all
three complex-structure deformations are polynomial. At
$(v,a,b)=(1/2,1/100,1/10)$ the face polynomials are
nondegenerate, including on the full torus; the exact check
tests their saturated critical ideals. Nondegeneracy is open,
so a maximal projective crepant partial toric resolution gives
a smooth projective Calabi--Yau threefold generically
\cite[Section~4.5]{Batyrev:1994}.
The small resolution constructed above is also smooth projective
and Calabi--Yau on the generic four-node family. The two models
are birational; hence they are connected by flops
\cite[Theorem~1]{Kawamata:2007Flops}.
\end{proof}

\paragraph{The fundamental period.}
An independent all-orders comparison is available. The three
nonnegative primitive relations among
\eqref{eq:x9-smooth-mirror-rays} are
\begin{equation}
 \begin{pmatrix}
 1&1&0&0&4&1&1\\
 0&0&1&0&2&1&0\\
 0&0&0&1&1&1&0
 \end{pmatrix}.
 \label{eq:x9-mutated-relations}
\end{equation}
They generate every nonnegative integral relation: its coefficients
satisfy $n_0=n_1=n_6=\alpha$, $n_2=\beta$, $n_3=\gamma$,
$n_4=4\alpha+2\beta+\gamma$ and $n_5=\alpha+\beta+\gamma$.
With the constant
term normalized as in \eqref{eq:x9-mutated-mirror}, take
\begin{equation}
 (\xi_1,\xi_2,\xi_3)=(va^2b^4,ab^2,ab),\qquad
 (v,a,b)=(\xi_1/\xi_2^2,\xi_3^2/\xi_2,\xi_2/\xi_3).
 \label{eq:x9-mutated-parameters}
\end{equation}
The integral exponent matrix has determinant $-1$.
The normalized constant-term period is
\begin{equation}
 W_{\rm sm}=\sum_{\alpha,\beta,\gamma\geq0}
 \frac{(8\alpha+4\beta+3\gamma)!\,
       \xi_1^\alpha\xi_2^\beta\xi_3^\gamma}
 {(\alpha!)^3\beta!\gamma!(4\alpha+2\beta+\gamma)!
                          (\alpha+\beta+\gamma)!}.
 \label{eq:x9-mutated-period}
\end{equation}
The substitution
$(\ell,m,n)=(\alpha,2\alpha+\beta+\gamma,4\alpha+2\beta+\gamma)$
identifies every coefficient with
\eqref{eq:x9-bulk-restricted-fundamental}. Its inverse is
$(\alpha,\beta,\gamma)=(\ell,n-m-2\ell,2m-n)$; nonnegativity
is exactly the support condition in that formula. Therefore
$W_{\rm sm}=W_0|_\Sigma$ to all orders, and analytically in their
common small-bulk domain of convergence.
